% source: An algebraic approach to the Riemann-Roch theorem and the arithmetic theory of function fields (Athanasios Papaioannou), http://www.fen.bilkent.edu.tr/~franz/mat/Riemann-Roch.pdf
% pdf_page: 0007
% retrieved: 2026-07-14
% transcribed_by: page-transcriber (AI vision, no OCR)

\documentclass[11pt]{article}
\usepackage{amsmath,amssymb,amsthm}

% Small-caps theorem style matching the source (DEFINITION headings, italic bodies).
\newtheoremstyle{scplain}{\topsep}{\topsep}{\itshape}{0pt}{\scshape}{.}{ }{\thmname{#1} \thmnote{#3}}
\theoremstyle{scplain}
\newtheorem*{definition}{Definition}

% Small-caps "Proof." to match the source (standard amsthm definition with \scshape head).
\makeatletter
\renewenvironment{proof}[1][\proofname]{%
  \par\pushQED{\qed}\normalfont\topsep6\p@\@plus6\p@\relax
  \trivlist\item[\hskip\labelsep\scshape #1\@addpunct{.}]\ignorespaces}%
  {\popQED\endtrivlist\@endpefalse}
\makeatother
\renewcommand{\qedsymbol}{$\square$}

% Non-standard operators used on this page (typeset upright in the source):
\DeclareMathOperator{\Div}{Div}
\DeclareMathOperator{\Princ}{Princ}
\DeclareMathOperator{\Class}{Class}
% Standard operators used: \deg, \dim, \sup. Standard amssymb symbols: \mathbb, \mathbf.

\begin{document}

\noindent
We may let $v_P : \Div(\mathbb{F}) \longrightarrow \mathbb{Z}$ be defined by
$v_P(D) = n_P$, the $P$'th coefficient of $D$. A partial ordering $\leq$ on
$\Div(\mathbb{F})$ can be introduced by writing $D_1 \leq D_2$ if and only if
$v_P(D_1) \leq v_P(D_2)$ for all $P \in \mathbf{P}$. A divisor $D$ is called
\emph{effective} if it is non-negative, in the sense that $D \geq 0$. Since the
formal sums are finite, we have a notion of \emph{degree} for divisors: if
\[
  D = \sum_{P \in \mathbf{P}_{\mathbb{F}}} n_P P,
\]
then we may define
\[
  \deg(D) = \sum_{p \in \mathbf{P}_{\mathbb{F}}} v_P(D) \deg(P).
\]
The degree map is actually a homomorphism $\Div(\mathbb{F}) \longrightarrow \mathbb{Z}$.

\begin{definition}[1.7]
Let $x \in \mathbb{F}^*$ and denote by $Z$ the set of zeroes of $x$ in
$\mathbf{P}_{\mathbb{F}}$. Then, we define the zero divisor of $x$ to be
\[
  (x)_0 = \sum_{P \in Z} v_P(x) P \in \Div(\mathbb{F});
\]
analogously, if $N$is the set of poles of $x$, we define the pole divisor of $x$
as $(x)_\infty = \sum_{P \in N}(-v_P(x))P$. Finally, the principal divisor of $x$
is defined to be
\[
  (x) = (x)_0 - (x)_\infty = \sum_{P \in \mathbf{P}_{\mathbb{F}}} v_P(x) P.
\]
\end{definition}

IN what follows, we will he assumption that $K$ is algebraically closed will now
imply that a non-zero element $x \in \mathbb{F}$ is constant if and only if
$(x) = 0$; this follows from the fact that any $z \in \mathbb{F}$ which is
transcendental over $\mathbb{K}$ has at least one zero and one pole and
Theorem 1.1.

\begin{definition}[1.8]
The subgroup $\Princ(\mathbb{F}) = \{(x) \mid x \in \mathbb{F}^*\}$ of
$\Div(\mathbb{F})$ is called the group of principal divisors of
$\mathbb{F}/\mathbb{K}$; the operation in this group is given by
$(xy) = (x) + (y)$. The class group
\[
  \Class(\mathbb{F}) = \Div(\mathbb{F})/\Princ(\mathbb{F})
\]
is called the divisor class group; its elements are classes of the form $[D]$,
where $D \in \Div(\mathbb{F})$ is a representative. Equality of divisor classes
$[D] = [D']$ in $\Class(\mathbb{F})$ induces an equivalence relation $D \sim D'$ in
$\Div(\mathbb{F})$. Therefore, $D \sim D'$ if and only if the difference $D - D'$
is some principal divisor $(x)$.
\end{definition}

Another object we need to introduce is the following class of $K$-vector spaces.

\begin{definition}[1.9]
For a divisor $D \in \Div(\mathbb{F})$, we let
\[
  L(A) = \{x \in \mathbb{F}^* \mid (x) \geq -D\} \cup \{0\}.
\]
\end{definition}

\noindent
\emph{Explicitly, if $v_P(D) > 0$ when and only when $P \in I$, where $I$ is a
finite subset of $\mathbf{P}_{\mathbb{F}}$ and $v_Q(D) < 0$ when and only when
$Q \in J$, where $J$ is a finite subset of $\Div(\mathbb{F})$, then $L(D)$ contains
all $x \in \mathbb{F}^*$ such that $x$ has zeroes of order at least $-v_Q(D)$ at
all elements $Q \in J$ and poles of order at most $v_Q(D)$ at all elements
$P \in I$.}

It's now obvious that $L(D)$ is a $K$-vector space, and all equivalent divisors
have isomorphic corresponding vector spaces (for this second assertion, assume
that $D \sim D'$ and then let $\phi : L(D) \longrightarrow L(D')$ be given by
$x \mapsto xz$, where $z$ is the fixed element of $\mathbb{F}$ such that
$D = D' + (z)$).

Our next objective will be to show that $\dim_{\mathbb{K}} L(D)$ is finite for all
divisors $D$ of $\mathbb{F}/\mathbb{K}$. This is intimately connected with the fact
that the \textbf{genus} of $\mathbb{F}/\mathbb{K}$, is in fact well-defined and
finite.

\begin{definition}[1.10]
The genus $g$ of a function field $\mathbb{F}/\mathbb{K}$ is defined by
\[
  g = \sup\{\deg D - \dim D + 1 \mid D \in \Div(\mathbb{F})\}.
\]
\end{definition}

\noindent
\emph{Note that this definition is legitimate (the supremum in question exists)
since the difference $\deg D - \dim D$ is bounded above, by the previous theorem.}

\begin{center}7\end{center}

\end{document}
